\documentclass{article}
\usepackage{graphicx} % Required for inserting images
\usepackage{amssymb} % more symbols
\usepackage{amsmath} %Required for implies and other symbols
\usepackage{amsthm} %Required for environments
\usepackage{tikz} %Diagrams
\usepackage{tikz-3dplot} %3D diagrams

\title{Notes on Linear Algebra Done Right! by S. Axler}
\author{Neilan Ghosh}
\date{June 2026}

\begin{document}
\maketitle
\section{Chapter 1 - Vector Spaces}
\subsection{Some basic properties (* is an arbitrary operation)} \\
1. Commutativity - $\alpha * \beta = \beta * \alpha$ \\
2. Associativity - $(\alpha*\beta)*\gamma = \alpha*(\beta*\gamma)$ \\
3. Identities - $\alpha * k = \alpha$, where k is a constant \\
4. Inverse - $\forall \alpha \in \mathbb{C}, \exists \beta \in \mathbb{C}$ such that $\alpha + \beta = 0$ \\
5. Distributive - $\lambda(\alpha*\beta) = \lambda\alpha * \lambda\beta$ \\
\subsection{Lists}
The set $\mathbb{R}^{2}$, which is a standard plane can be set up as \\
\begin{equation}
    \mathbb{R}^2 = \{(x,y): x,y \in \mathbb{R}\}
\end{equation}

Similarly, the set $\mathbb{R}^3$, which is ordinary 3D-space can be set up as
\begin{equation}
    \mathbb{R}^3 = \{(x,y,z): x,y,z \in \mathbb{R}\}
\end{equation}

A list is an ordered collection of n elements \\ 
The "coordinates" in n-dimensional space is called a list. You can think of it as "arguments for movement". \\
\begin{center}
    List of n-length tuple $\rightarrow (z_1, z_2, z_3,...z_n)$ = z 
\end{center}
\begin{itemize}
    \item Sometimes, lists are contracted to a single letter as shown above 
    \item A list of length 0 is (). This does not lead to trivial exceptions in some theorems. 
\end{itemize}

List v Sets
\begin{itemize}
    \item (3,5) $\neq$ (5,3), but \{3,5\} = \{5,3\}
    \item (1,1) $\neq$ (1,1,1), but \{1,1\} = \{1,1,1\} = \{1\}
\end{itemize}

\subsection{$\mathbb{F}^n$}
Use $\mathbb{F}$ whenever $\mathbb{R}$ or $\mathbb{C}$ ($\mathbb{F}$ implies whichever space is more appropriate). $\mathbb{F}^n$ is the set of all n-tuple lists with elements belonging to $\mathbb{F}$ 
\begin{center}
    $\mathbb{F}^n$ = \{($x_1,x_2,x_3,...x_n$): $x_k$ $\in \mathbb{F}$ for k = 1,2,3,...n\}
\end{center}

We say that $x_k$ is the kth coordinate of the list for ($x_1,x_2,...x_n$) $\in$ $\mathbb{F}$ and k $\in \{1,2,3...n\}$ \\ 
Even if our visualization and geometric approach stops at $\mathbb{R}^4$ and $\mathbb{C}^2$, we can still use algebraic manipulations. \\ 
For x $\in \mathbb{F}^n$, the additive inverse of x is the vector -x $\in \mathbb{F}^n$ such that 
\begin{center}
    x + (-x) = 0
    $\therefore$ If x = $(x_1, x_2,..x_n)$ $\implies$ $-x = (-x_1,-x_2,...-x_n)$
\end{center}
Product of a number $\lambda$ and a vector in $\mathbb{F}^n$ is done by multiplying individually. 
\subsection{What is a field?}
\textbf{Definition} \\ 
\begin{center}
    "A set containing at least two distinct elements, 0 and 1, along with operations of addition and multiplication following the laws"
\end{center}
\textbf{Intuition:}
Think of it as a counter-top. Your ingredients (elements) can be combined (in ways called operations) to give new ingredients which can be used further. The counter-top with ingredients and combining methods is the field. \\
Logically, you cannot get something out of the set you are operating out of. \\
A reason why $\mathbb{Z}$ is not a field is because most elements of it do not have a multiplicative inverse which belongs to $\mathbb{Z}$ itself. \\
Also, remember that there is an "undo" button for all ingredients except 0 (which are the additive and multiplicative inverses) \\

\subsection{Vector Spaces}
\textbf{Incentive}: Properties of addition and multiplication ($\because$ Commutativity, associativity, existence of identity and inverse, distributivity) 
\begin{center}
    A set V with operations of addition and scalar multiplication, which follow the rules of associativity, commutativity etc.
\end{center}
\textbf{What this means} is that I can add two elements and multiply them with an element of $\mathbb{F}$ (Scalar, can be $\mathbb{R}$ or $\mathbb{C}$) \\ 
Mathematically, I would say \\
\begin{itemize}
    \item Addition is a function that gives an element t (can be unique or equal) $\in$ V such that t = u + v where u,v $\in$ V
    \item Scalar multiplication is a function that gives an element to $\lambda$v $\in$ V s.t. v $\in$ V and $\lambda\in \mathbb{F}$ 
\end{itemize}
\textbf{Definition:} 
\begin{center}
    A vector space V is a set along with an addition on V and a scalar multiplication s.t. the following hold \\
    \textbf{Commutativity} - $u+v = v+u, \forall u,v \in V$ \\
    \textbf{Associativity} - $(u+v)+w = u + (v+w)$ and $(ab)v = a(bv) \forall u,v,w \in V$ and $\forall a,b \in \mathbb{F}$ \\
\end{center}
Along with this, it must have additive identities, additive inverses, multiplicative identities and the distributive property

\subsection{Isn't this just a field?}
\textbf{NO.} A few reasons for this are 
\begin{itemize}
    \item In a field, the elements can be added and multiplied i.e. $ab \in \mathbb{F} \forall a,b\in \mathbb{F}$. But in a vector space, two vectors CANNOT be multiplied together i.e. vw is not defined $\forall v,w \in V$
    \item You might think "Why not define multiplication as cross product? Now, we can multiply vectors and get a vector out of it." \\
    NO. This does not work since cross products are defined in 3D or 7D but not in other dimensions. On top of that, the associativity is violated in the vector space V.
\end{itemize}
Also, whenever we say "V is a vector space over $\mathbb{F}$", it implies that for scalar multiplication, we take elements from the field $\mathbb{F}$
\subsection{An interesting example of a vector space}
A set of functions $\rightarrow \mathbb{F}^S$ \\
If S is a set, then $\mathbb{F^S}$ is the set of functions $S \rightarrow \mathbb{F}$ \\
ADDITION $\rightarrow \forall f,g\in \mathbb{F}^S$, then f+g is defined as 
\begin{center}
    (f+g)(x) = f(x) + g(x) $\forall x \in S$
\end{center}
SCALAR MULTIPLICATION $\rightarrow \forall f\in \mathbb{F}^S$ and $\lambda f \in \mathbb{F}^S$ is defined as 
\begin{center}
    $(\lambda f)(x)$ = $\lambda f(x) \forall x\in S$
\end{center}
\subsection{Excerpt from Axler}
\begin{center}
    "The vector space $\mathbb{F}^n$ is a special case of the vector space $\mathbb{F}^S$ because each $(x_1, ..., x_n) \in \mathbb{F}^n$ can be thought of as a function x from the set \{1,2,...,n\} to F by writing x(k) instead of $x_k$ for the kth coordinate of $(x_1, ..., x_n)$."
\end{center}
He is trying to tell you to move on from "arrows" to functions. \\
$\mathbb{F}^n$ is a set of functions from \{1,2,3,...n\} to $\mathbb{F}$ \\
Say x:\{1,2,3,...n\} $\rightarrow \mathbb{F}$ \\
x(1) = $x_1$, x(2) = $x_2$,... \\ 
This swaps the coordinate logic with functions. \\
\textbf{Incentive:} Used when there is no definite coordinate system in place. \\ 
\\
When we defined the vector space V, we set two operations in the space itself: Addition and scalar multiplication \\

\textbf{Claim:} Prove that a vector space has a unique additive identity \\
\textbf{Proof:} Assume that 0 and 0' are additive identities for some vector space V. \\
Then, 
\begin{center}
    0' = 0+0' = 0'+0 = 0
\end{center}
The first equality holds since 0 is an additive identity. \\
The second equality holds due to commutativity. \\
The third equality holds since 0' is an additive identity.\\
$\therefore$ There exists only one additive identity for a vector space. \\
\\
\textbf{Claim:} P.T. every element in a vector space has a unique additive inverse. \\
\textbf{Proof:} Assume that w and w' $\in V$ are the additive inverses of an element v $\in V$ \\
$\implies$ w = w+0 = w+(v+w') = (w+v)+w' = 0+w' = w' \\ 
$\therefore$ w = w' \\
$\therefore$ There exists a unique additive inverse for every element in a vector space \\
\textit{Why not just say v+w = v+w'? Subtracting on both sides, we get w = w'} \\
Subtraction has not yet been defined. Same with inverses (don't add by inverses). \\ 
\textbf{Consequence of proving additive inverse unique:} We can now say for v,w $\in V$ 
\begin{center}
    -v denotes additive inverse of v \\
    w-v is defined to be w+(-v)
\end{center}
\textbf{Claim:} P.T. 0*v = 0 $\forall v \in V$ \\
\textbf{Proof:} For $\forall v \in V$, we have \\
\begin{center}
    0v = (0+0)v = 0v+0v
\end{center}
Adding the additive inverse of 0v on both sides, we get \\
\begin{center}
    $\implies$ 0 = 0v \\
    Proved.
\end{center}
\textbf{Claim:} P.T. (-1)v = -v $\forall v \in V$ \\
\textbf{Proof:} For $v \in V$, we know that 
\begin{center}
    (-1)v + v = (-1)v + 1v = ((-1) + 1)v = 0v = 0 \\
    $\implies$ (-1)v + v = 0 \\ 
    Adding additive inverse of v on both sides, \\ 
    $\implies$ (-1)v = -v \\
    Proved.
\end{center}
\subsection{Subspaces}
\textbf{Definition:}
\begin{center}
    A subset U of V is called a subspace of V if U is also a vector space with the same additive identity, addition and scalar multiplication as on V.
\end{center}
\textbf{Intuition:} \\
Remember the counter-top? Now, you segment it; some of it is meant for cutting, some for spices etc. You can still do the same things (combine as in operations) but now in a specific region which belongs to the larger counter-top (the vector space $V$) \\
\textbf{Conditions:} \\
A subset U of V is a subspace $\iff$ U satisfies the following 
\begin{itemize}
    \item Additive Identity \\
    0 $\in U$            ...(1)
    \item Closed under addition \\
    u,w $\in U \rightarrow$ u + w $\in U$             ...(2)
    \item Closed under multiplication \\ 
    u $\in U$, a $\in \mathbb{F} \rightarrow$ au $\in U$        ....(3)
\end{itemize}
(1) signifies the additive identity is present in $U$ \\
(2) signifies that addition makes sense \\
(3) signifies that multiplication makes sense \\ 
Notice that, since we proved associativity, commutativity etc. for vector space V, any subspace U of V would automatically have it. Whatever you proved recently applies to the subspace U. \\
Ex: $\forall$ u  $\in U \rightarrow$ (-1)u = -u $\in U$ \\
\subsection{Examples of Subspaces}
\textbf{1.} If b $\in \mathbb{F}$, then 
\begin{center}
    $\{(x_1, x_2,x_3,x_4) \in \mathbb{F}^4 : x_3 = 5x_4 + b\}$ \\
    is a subspace of $\mathbb{F}^4 \iff$ b = 0
\end{center}
\textbf{Note:} Algebraically, I would say \\
\textbf{Case 1:} b = 0 $\implies x_3 = 5x_4$ \\
(0,0,0,0) does belong to this subspace as 0 = 5(0) \\
\textbf{Case 2:} b $\neq$ 0 (say b = 1) $\implies x_3 = 5x_4 +1$  \\
$\therefore$ (0,0,0,0) does not belong to this subset, disqualifying it as a subspace. \\ 
\textit{Geometrically, the origin being there is imperative to make a subset a subspace.} \\
\\
\textbf{2.} Set of continuous real valued functions on the interval [0,1] is a subspace of $\mathbb{R}^{[0,1]}$ \\
\\
\textbf{3.} Set of differentiable real-valued functions on $\mathbb{R}$ is a subspace of $\mathbb{R}^{\mathbb{R}}$
\subsection{Geometric Interpretation of subspaces}
The subspaces of $\mathbb{R}^2$ are any lines which pass through (0,0), $\{0\}$ itself, and $\mathbb{R}^2$ \\
Note that other shapes such as parabolae, hyperbolas etc. DO NOT come under subspace classification. \\

\textbf{Question:} Is S = $\{(x,y): x \ge 0, x,y \in \mathbb{R}\}$ a subspace of $\mathbb{R}^2$? \\
NO. Go through it. Origin (0,0) is there. Addition makes sense (is closed) here, but scalar multiplication falls apart. (Multiply by -1 and see) 
\subsection{Sums of Subspaces}
Generally, we are interested in subspaces of a vector space. 
\textbf{Definition:} 
\begin{center}
    Suppose $V_1,..., V_m$ are subspaces of V. The sums of $V_1,...V_m$ denoted by $V_1+...+ V_m$ is the set of all possible sums of elements of $V_1,..., V_m$. More precisely, \\
    $V_1+...+ V_m$ = $\{v_1+...+ v_m:v_1 \in V_1,...v_m \in V_m \}$
\end{center}
\textbf{Note:} You would think of sum of subspaces to be the union of the subsets. Do not do this mistake. It holds a different logic. \\
Union is "Either A or B" \\
Sum of subspaces is "Move in A then B" \\
\\
Example: \\
1. Sum of subspaces of $\mathbb{F}^3$ \\
\begin{center}
    U = $\{(x,0,0) \in \mathbb{F}^3 : x \in \mathbb{F}\}$ \\
    and V = $\{(0,y,0) \in \mathbb{F}^3 : y \in \mathbb{F}\}$ \\ 
    $\implies$U + W = $\{(x,y,0) \in \mathbb{F}^3 : x,y \in \mathbb{F}\}$
\end{center}    
2. A sum of subspaces of $\mathbb{F}^4$ 
\begin{center}
    U = $\{(x,x,y,y) \in \mathbb{F}^4: x,y \in \mathbb{F}\}$ and V = $\{(x,x,x,y) \in \mathbb{F}^4 : x,y \in \mathbb{F}\}$ 
\end{center}
You cannot directly say that 
\begin{center}
    U + W = $\{(2x,2x,x+y,2y):\mathbb{F}^4 : x,y \in \mathbb{F}\}$        ...(1)
\end{center}
Axler points out that 
\begin{center}
    An element of U is (a,a,b,b) \\
    An element of W is (c,c,c,d) \\
    $\therefore$ An element of U+W is ((a+c),(a+c),(b+c),(b+d)) \\
    $\therefore$ U+W $\subseteq$ $\{(x,x,y,z)\in \mathbb{F}^4: x,y,z \in \mathbb{F}\}$
\end{center}
To prove inclusion in the other direction, $x,y,z \in \mathbb{F}$ 
\begin{center}
    (x,x,y,z) = (x,x,y,y) + (0,0,0,(z-y))
\end{center}
The second term $\in$ Subspace W 
\begin{center}
    $\therefore$ U + W = $\{(x,x,y,z) \in \mathbb{F}^4 : x,y,z \in \mathbb{F}\}$
\end{center}
Extremely important to know that when you go fast, you get (1), WHICH IS WRONG. \\
Go with set equality by double inclusion. \\
\\
\textbf{Claim:} Supposed $V_1,...V_m$ are the subspaces of V. Then, P.T. $V_1+...+V_m$ is the smallest subspace containing $V_1,...V_m$. \\
\textbf{Note:} This proof will be slightly informal to help understanding. \\
\begin{center}
    \tdplotsetmaincoords{70}{120}
\begin{tikzpicture}[tdplot_main_coords]

\draw[->] (0,0,0)--(3,0,0) node[right] {x};
\draw[->] (0,0,0)--(0,3,0) node[left] {y};
\draw[->] (0,0,0)--(0,0,1) node[above] {z};

\filldraw[blue!15,opacity=0.7]
(0,0,0)--(0,2.5,0)--(2.5,2.5,0)--(2.5,0,0)--cycle;
\node[blue] at (1.2,1.2,0) {$V_1+V_2$};

\draw[red,thick]
(0,0,0)--(2.5,0,0);
\node[red] at (2,0,0.3) {$V_1$};

\draw[green,thick]
(0,0,0)--(0,2.5,0);
\node[green] at (0,2,0.3) {$V_2$};

\end{tikzpicture}
\end{center}
We will just be operating with $V_1,V_2$ and $V_1+V_2$ where V = $\mathbb{R}^3$ \\
First of all, we can comfortable say $V_1,V_2$ and $V_1+V_2$ are subspaces of V (Additive identity,closed addition and scalar multiplication)\\
Secondly, $V_1$ and $V_2$ are contained in $V_1+V_2$ (You can prove this algebraically also). \\
Now, conversely, every subspace of V containing $V_1,V_2$ must contain $V_1+V_2$ (The blue plane). \\
No matter what subspace you take, the smallest one containing $V_1,V_2$ must be $V_1+V_2$. \\
Think about it. Any other subspace I take would contain $V_1+V_2$ naturally. The largest subspace of V I can take is V itself. The smallest subspace of V containing $V_1,V_2$ must be $V_1+V_2$. 
\subsection{Direct Sums}
\textbf{Definition:}
\begin{center}
    Suppose $V_1,...V_m$ are subspaces of V
    \begin{itemize}
        \item The sum $V_1+...+V_m$ is called a direct sum if each element of $V_1+...+V_m$ can be written in only one way as a sum $v_1,...v_m$, where each $v_k \in V_k$.
        \item If $V_1+...+V_m$ is a direct sum, then $V_1\oplus...\oplus V_m$ denotes $V_1+...+V_m$ with $\oplus$ saying that this is a direct sum.
    \end{itemize}
\end{center}
\textbf{Logic:} \\
Why do we need this? The reason is uniqueness. \\
Without uniqueness, 
\begin{center}
    Suppose V = U + W \\
    $\implies$ v = u + w, where $v \in V$, $u \in U$, $w \in W$ 
\end{center}
What can also happen is 
\begin{center}
    v = u + w = u' + w'
\end{center}
For example, 
\begin{center}
    (1,2) = (1,0) + (0,2) = (10,5) + (-9,-3) = ...
\end{center}
With uniqueness, 
\begin{center}
    i.e. when V = U $\oplus$ W \\
    only v = u + w is possible 
\end{center}
This is what the first point is about. 
\subsection{Examples of Direct Sums}
1. U = $\{(x,y,0)\in \mathbb{F}^3: x,y \in \mathbb{F}\}$ and W = $\{(0,0,z) \in \mathbb{F}^3: z \in \mathbb{F}\}$ \\
Clearly, $\mathbb{F}^3 = U \oplus W$ \\
2. Suppose $V_k$ is a subspace of $\mathbb{F}^n$ such that the only non-zero value is the kth element of the list. \\
\begin{center}
    i.e. $V_2$ = (0,y,0,0,...,0) \\
    $\therefore \mathbb{F}^n = V_1 \oplus V_2 \oplus V_3 \oplus ... \oplus V_n$ 
\end{center}
\subsection{A non-example of Direct Sums}
1. Suppose
\begin{center}
    $V_1$ = $\{(x,y,0)\in \mathbb{F}^3: x,y \in \mathbb{F}\}$ \\
    $V_2$ = $\{(0,0,z)\in \mathbb{F}^3: z \in \mathbb{F}\}$ \\
    $V_3$ = $\{(0,y,y)\in \mathbb{F}^3: y \in \mathbb{F}\}$
\end{center}
Already, we can see that $\mathbb{F}^3 = V_1 \oplus V_2$ \\
$V_3$ only add a possibility of non-uniqueness.
\begin{center}
    (0,0,0) = (0,1,0) + (0,0,1) + (0,-1,-1)
            = (0,0,0) + (0,0,0) + (0,0,0)
\end{center}
It gives a way to modify $V_1 + V_2$ which allows the non-uniqueness due to the overlap. \\
\\
\textbf{Claim:} Suppose $V_1,...V_m$ are subspaces of V. P.T. $V_1+...+V_m$ is a direct sum $\iff$ the only way to write 0 as a sum $V_1+...+V_m$ where each $v_k \in V_k$ is by taking $v_k$ = 0 \\
\textbf{Proof:} Left to right is fairly obvious. 
\begin{center}
    In order for $V_1+...+V_m$ to be a direct sum, there must only be one way (trivial) to write 0 as $V_1+...+V_m$ \\
\end{center}
Right to left: \\
Let v $\in V$ such that 
\begin{center}
    v = $v_1+...+v_m$   ...(1) \\
    where $v_1 \in V_1, v_2 \in V_2$ ...
\end{center}
Assume that, for $u_1 \in V_1$, $u_2 \in V_2$, ...
\begin{center}
    v = $u_1+...+u_m$
\end{center}
Subtracting (1) and (2), \\
\begin{center}
    0 = $(v_1-u_1) + (v_2-u_2)+...(v_m-u_m)$
\end{center}
The only way to do this is trivially i.e. 
\begin{center}
    $v_k - u_k = 0$ for k =1,2,3,...
    $\therefore v_k = u_k$ \\
    Proved.
\end{center}
\textbf{Remark/Reason:} Subtraction is actually a common way to generate such proofs. Simply put, we do it here to get the 0 and get the condition of direct sums. Simple operation, and easy to miss also. \\
\\
\textbf{Claim:} Suppose U and W are subspaces of V.\\
Then, P.T. U + W is a direct sum $\iff$ $U \cap W = \{0\}$ \\
\textbf{Proof:} Left to right: \\
Assume that U + W is a direct sum. Let v $\in U \implies -v \in W$
\begin{center}
    v + (-v) = 0 \\
    ($\because v \in U \cap W \implies v \in U \implies v \in W \implies -v \in W$) \\
    $\implies$ u + w = 0 \\
    Proved.
\end{center}
Right to left: 
\begin{center}
    Assume that $U \cap W = \{0\}$ \\
    0 = u + w $\forall u \in U, w \in W$ \\
    $\implies$ u = -w \\
    $\implies$ -u = w $\in W$
    $\implies$  u $\in W \cap U$ \\
    $\therefore$ u = 0 = w \\
    Proved.
\end{center}
\textbf{Note:} Read this properly. It is a good proof and also teaches you about iff proofs. 
\subsection{A very interesting point by Axler}
\begin{center}
    Sums of subspaces are analogous to unions of subsets. Direct sums are analogous to disjoint unions of subsets. 
\end{center}
Also, when dealing with more than two subspaces, we cannot use the same logic as we did earlier. It is slightly different (You will see this later).



\end{document}